\subsubsection{SEM Implementation}\label{SEC:SEM:implementation}

This section provides details on certain aspects of the synthetic eddy method
regarding how it is implemented in GFR.
%

%###############################################################################

\subsubsubsection{Interpolation of SEM Profile onto Inflow Plane}%
[SEC:SEM:implementation:profile]

GFR requires a 1D The SEM profile 

%###############################################################################

\subsubsubsection{Inflow Plane Geometry}[SEC:SEM:implementation:geometry]

In its current form in GFR, the SEM can only used on a constant-$x$ inflow
boundary.
%
The primary reason for this restriction is that the synthetic eddies are
currently only propagated in $x$.

%%%%%%%%%%%%%%%%%%%%%%%%%%%%%%%%%%%%%%%%%%%%%%%%%%%%%%%%%%%%%%%%%%%%%%%%%%%%%%%%

Another geometry restriction for using the SEM is that it is highly recommended
that the inflow boundary be planar; this restriction will almost certainly be a
permanent requirement even if the constant-$x$ restriction is removed.
%
The algorithm that interpolates the SEM profile onto the inflow boundary points
uses the wall distance and wall normal to rotate and interpolate the profile.
%
The wall distance part of this algorithm could potentially cause problems with a
non-planar inflow boundary.
%
For example, imagine an inflow boundary that is highly convex from the interior
domain.
%
The wall distance for the boundary points will be larger than expected, so the
interpolated SEM profile might not be correct.
%
Again, this is a recommendation, so it might work, but there will probably be
some inaccuracies or unexpected consequences.

%###############################################################################

\subsubsubsection{Efficiency of SEM}[SEC:SEM:implementation:efficiency]

Because the eddies are propagated linearly through the SEM box, as soon as a new
eddy is created, it can be determined which face points will be affected by the
eddy.
%
In a crude implementation of the SEM algorithm, the fluctuating velocities at
each boundary point are computed using an outer loop over each boundary point
and an inner loop over each eddy to evaluate each eddy's contribution to the
fluctuations at the current boundary point.

%###############################################################################

\subsubsubsection{Parallel Details}[SEC:SEM:implementation:parallel]

The SEM was chosen because it is simple and efficient, and its point-based
formulation allows for it to be used in parallel.
%
However, care needs to be taken if the cells on the SEM boundary are split
across multiple partitions in order to prevent significant communication or
synchronization between the grid partitions sharing the inflow boundary.
%
Each synthetic eddy has a zone of influence, and any SEM boundary face point
within this zone must include the contribution of this random eddy into its
summation of all eddies affecting the face point.
%
\Fref{FIG:SEM:ParallelSEM} shows a possible example of the zone of influence of
a synthetic eddy affecting face points of cells from different partitions.
%
\begin{figure}[tb]
  \centering
  \insertimage[width=0.75\textwidth]{tikz/sem_partition_boundary}
 %\documentclass[border=0pt]{standalone}

\usepackage{tikz}        % drawing macros
  \usetikzlibrary{math}
  \usetikzlibrary{calc}  % calculations
  \usetikzlibrary{arrows}  % arrowheads
  \usetikzlibrary{arrows.meta}  % arrowheads
  \usetikzlibrary{decorations.markings}
\usepackage{tikz-3dplot} % drawing macros 3d
\usepackage{pgfplots}
%\pgfplotsset{compat=1.18}

\begin{document}

\tikzmath{
  \lencell = 2;
  \ybxval = 4;
  %
  \ybxmx = \ybxval;
  \ybxmn = -1 * \ybxmx;
  %
  \yfigmx = 0.2 + \ybxval;
  \yfigmn = -1 * \yfigmx;
  %
  \lgdtxt = 6.0;
  \lgdsym = \lgdtxt - 0.5;
  \lgdmnsym = \lgdsym - 0.2;
  \lgdplsym = \lgdsym + 0.2;
}
%
\begin{tikzpicture}
  [scale=0.9,
   cell/.style = {draw=black,thick},
  %sembox/.style = {draw=semboxblue,very thick,fill=semboxblue,semitransparent},
   sembox/.style = {draw=semboxblue,very thick,fill=semboxblue,nearly opaque},
  %sembox/.style = {draw=semboxblue,very thick,fill=semboxblue,very nearly opaque},
  %sembox/.style = {draw=semboxblue,very thick,fill=semboxblue,ultra nearly opaque},
   sembnd/.style = {draw=sembndblue,thick,fill=sembndblue},
   clprta/.style = {draw=partitiona,fill=partitiona,thick},
   clprtb/.style = {draw=partitionb,fill=partitionb,thick},
   facpnt/.style = {rectangle,fill=black,minimum size=1.5mm,inner sep=0pt},
   eddypt/.style = {circle,fill=eddyred,minimum size=1mm,inner sep=0pt},
   zoninf/.style = {circle,draw=zonelinered,thick,fill=zonered,minimum size=10mm,inner sep=0pt},
   legend/.style = {font=\normalsize,anchor=west},
   refpnt/.style = {minimum size=1mm,inner sep=0pt}]

  %
  \definecolor{yellowfill}{rgb}{0.929411765,0.921568627,0.878431373}
  \definecolor{yellowdraw}{rgb}{0.768627451,0.741176471,0.592156863}
 %\definecolor{semboxblue}{rgb}{0.815686274,0.901960784,1.0}
 %\definecolor{semboxblue}{rgb}{0.815686274,0.952941176,1.0}
  \definecolor{semboxblue}{RGB}{208,243,255}
  \definecolor{sembndblue}{RGB}{0,103,206}
  \definecolor{zonered}{RGB}{255,185,185}
  \definecolor{zonelinered}{RGB}{255,141,141}
  \definecolor{eddyred}{RGB}{175,0,0}
  \definecolor{partitiona}{RGB}{204,255,204}
  \definecolor{partitionb}{RGB}{255,255,204}
  %
  \useasboundingbox (-1.8,\yfigmn) rectangle (14,\yfigmx);
  %
  \draw [clprta] ( 0, 0) rectangle (4,4);
  \draw [clprtb] ( 0, 0) rectangle (4,-4);
  \draw [sembox] (-1.6,-4) rectangle (1.6,4);
  %
  \node at (-0.20000,0      ) [zoninf] {};
  \node at (-1.21512,-1.2037) [zoninf] {};
  \node at ( 0.09224, 1.1088) [zoninf] {};
  \node at (-1.52784, 2.4137) [zoninf] {};
  \node at ( 1.47928,-1.8615) [zoninf] {};
  \node at ( 0.20160,-2.4445) [zoninf] {};
  \node at ( 0.74632, 2.9637) [zoninf] {};
  \node at ( 1.21032, 1.7112) [zoninf] {};
  %
  \node at (-0.20000,0      ) [eddypt] {};
  \node at (-1.21512,-1.2037) [eddypt] {};
  \node at ( 0.09224, 1.1088) [eddypt] {};
  \node at (-1.52784, 2.4137) [eddypt] {};
  \node at ( 1.47928,-1.8615) [eddypt] {};
  \node at ( 0.20160,-2.4445) [eddypt] {};
  \node at ( 0.74632, 2.9637) [eddypt] {};
  \node at ( 1.21032, 1.7112) [eddypt] {};
  %
  \draw [sembnd] (-1.6,-4) -- (-1.6,4);
  \draw [sembnd] ( 1.6,-4) -- ( 1.6,4);
  %
  % Cell boundaries
  \draw [cell] ( 0, 0) -- ( 4, 0);
  \draw [cell] ( 0, 4) -- ( 0,-4);
  \draw [cell] ( 0,-3.2) -- ( 4,-3.2);
  \draw [cell] ( 0, 3.2) -- ( 4, 3.2);
  \draw [cell] ( 3.2,-4) -- ( 3.2, 4);
  %
  % Face points on SEM boundary
  \node at (0,-0.24) [facpnt] {};
  \node at (0, 0.24) [facpnt] {};
  %
  \node at (0,-2.96) [facpnt] {};
  \node at (0, 2.96) [facpnt] {};
  %
  \node at (0,-0.96) [facpnt] {};
  \node at (0, 0.96) [facpnt] {};
  %
  \node at (0,-2.24) [facpnt] {};
  \node at (0, 2.24) [facpnt] {};
  %
  \node at (0,-3.44) [facpnt] {};
  \node at (0, 3.44) [facpnt] {};
  %
 %\node at (-0.4,0.5)  [legend] {Cells on Partition A};
 %\node at (-0.4,-0.5) [legend] {Cells on Partition B};
  %
  \draw [sembox] (\lgdmnsym, 2.2) rectangle (\lgdplsym, 2.8);
  \draw [sembnd] (\lgdmnsym, 2.2) --        (\lgdmnsym, 2.8);
  \draw [sembnd] (\lgdplsym, 2.2) --        (\lgdplsym, 2.8);
  %
  \draw [clprta] (\lgdmnsym,-1.7) rectangle (\lgdplsym,-1.3);
  \draw [cell]   (\lgdmnsym,-1.7) rectangle (\lgdplsym,-1.3);
  %
  \draw [clprtb] (\lgdmnsym,-2.7) rectangle (\lgdplsym,-2.3);
  \draw [cell]   (\lgdmnsym,-2.7) rectangle (\lgdplsym,-2.3);
  %
  \node at (\lgdsym, 1.5) [scale=1.75,eddypt] {};
  \node at (\lgdsym, 0.5) [scale=0.5,zoninf] {};
  \node at (\lgdsym,-0.5) [scale=1.5,facpnt] {};
  %
  \node at (\lgdtxt, 2.5) [scale=1.0,legend] {SEM Virtual Box};
  \node at (\lgdtxt, 1.5) [scale=1.0,legend] {Synthetic eddy};
  \node at (\lgdtxt, 0.5) [scale=1.0,legend] {Zone of influence for eddy};
  \node at (\lgdtxt,-0.5) [scale=1.0,legend] {Face point of cell on SEM boundary};
  \node at (\lgdtxt,-1.5) [scale=1.0,legend] {Grid Cells from partition A on SEM boundary};
  \node at (\lgdtxt,-2.5) [scale=1.0,legend] {Grid Cells from partition B on SEM boundary};
  %
 %\draw [cell] (-1.8,\yfigmn) rectangle (14,\yfigmx);
  %
  % SEM box x-interval of [-3,-1]
%%\draw [sembox] (-3,-3) rectangle (-1,3);
%%%
%%\node at (-2.125   ,0     ) [zoninf] {};
%%\node at (-2.75945,-1.2037) [zoninf] {};
%%\node at (-1.94235, 1.1088) [zoninf] {};
%%\node at (-2.9549 , 2.4137) [zoninf] {};
%%\node at (-1.07545,-1.8615) [zoninf] {};
%%\node at (-1.874  ,-2.4445) [zoninf] {};
%%\node at (-1.53355, 2.9637) [zoninf] {};
%%\node at (-1.24355, 1.7112) [zoninf] {};
%%%
%%\node at (-2.125  ,0      ) [eddypt] {};
%%\node at (-2.75945,-1.2037) [eddypt] {};
%%\node at (-1.94235, 1.1088) [eddypt] {};
%%\node at (-2.9549 , 2.4137) [eddypt] {};
%%\node at (-1.07545,-1.8615) [eddypt] {};
%%\node at (-1.874  ,-2.4445) [eddypt] {};
%%\node at (-1.53355, 2.9637) [eddypt] {};
%%\node at (-1.24355, 1.7112) [eddypt] {};
%%%
%%\draw [sembnd] (-3,-3) -- (-3,3);
%%\draw [sembnd] (-1,-3) -- (-1,3);
%%%
%%\draw [cell] (-2,-3) -- (-2,3);
%%\draw [cell] (-2,0) -- (3,0);
%%%
%%% Face points on partition B
%%\node at (-2,-0.3) [facpnt] {};
%%\node at (-2,-1.4) [facpnt] {};
%%\node at (-2,-2.8) [facpnt] {};
%%%
%%% Face points on partition A
%%\node at (-2,0.3)  [facpnt] {};
%%\node at (-2,1.4)  [facpnt] {};
%%\node at (-2,2.8)  [facpnt] {};
%%%
%%\node at (-1.4,0.5)  [legend] {Cells on Partition A};
%%\node at (-1.4,-0.5) [legend] {Cells on Partition B};
%%%
%%\draw [sembox] (4.8,1.7) rectangle (5.2,2.3);
%%\draw [sembnd] (4.8,1.7) -- (4.8,2.3);
%%\draw [sembnd] (5.2,1.7) -- (5.2,2.3);
%%\draw [cell]   (4.8,-2.2) rectangle (5.2,-1.8);
  %
\end{tikzpicture}

\end{document}

  \caption{SEM box spanning multiple partitions.}
  \label{FIG:SEM:ParallelSEM}
\end{figure}

%%%%%%%%%%%%%%%%%%%%%%%%%%%%%%%%%%%%%%%%%%%%%%%%%%%%%%%%%%%%%%%%%%%%%%%%%%%%%%%%

Every time an eddy leaves the virtual box, a new eddy, with a random intensity,
is created at a random location on the leading face of the box.
%
These two random operations are dependent on a random number generator which is
initialized at the beginning of the simulation using a random seed that is
consistent across all partitions.
%
However, the consistency across all partitions of the random generation of
synthetic eddies will quickly breakdown if each partition only accounts for
eddies that affect its own cells.
%
Therefore, all partitions must maintain a common database of all the synthetic
eddies including the randomness of all eddies that are regenerated at the
leading face of the virtual SEM box.

%%%%%%%%%%%%%%%%%%%%%%%%%%%%%%%%%%%%%%%%%%%%%%%%%%%%%%%%%%%%%%%%%%%%%%%%%%%%%%%%

When computing the velocity fluctuations at each face point due to the eddies
influencing a given point, looping through the many thousands of eddies to find the
dozens or so that could actually influence the face point at a given time
becomes computationally expensive across thousands of face points on the SEM
boundary.
%
While all partitions must maintain their own consistent global database of all
the synthetic eddies in the simulation, each partition must maintain secondary
localized databases for each face point on the partition for just those eddies
whose zone of influence might affect the face point.
%
Since these eddies are constantly being regenerated at a random location after
propagating through the SEM box, each of these secondary localized databases
have to be constantly updated.
%
By using these secondary localized databases, each face point loops through only
those eddies that can have a possible influence on the face point when computing
the SEM velocity fluctuation, thus significantly improving the parallel
efficiency of the SEM method.
%
In testing this parallel implementation of the SEM method, we were able to use
hundreds of thousands of synthetic eddies across thousands of MPI processes with
minimal computational cost to the overall simulation.

%%%%%%%%%%%%%%%%%%%%%%%%%%%%%%%%%%%%%%%%%%%%%%%%%%%%%%%%%%%%%%%%%%%%%%%%%%%%%%%%

Tracking and communicating the interaction between thousands of eddies
and thousands of face points can quickly become computationally expensive.
%
Therefore, the parallel SEM implementation in GFR is designed for both
accuracy and efficiency.
%
When an inflow boundary is split across multiple partitions,
communication costs are eliminated by having each one store and track a
copy of the same virtual SEM volume, eddies, and random seed.
%
Because the eddies are convected in a uniform direction,
as each new eddy is created one can project it onto the
physical inflow boundary and know which face point(s) it will
eventually influence.
%
So, to reduce the search time between eddies and face points that
must be done every time step, each face point tracks its own
secondary database of eddies that may affect it.
%
%Eddies need only be added\slash removed from the secondary database only
%when they enter\slash exit the virtual volume.
%
In testing this parallel implementation of the SEM method,
hundreds of thousands of synthetic eddies were used
across thousands of MPI processes with minimal computational cost
to the overall simulation.

%===============================================================================

Note that even though the synthetic eddy method is designed to
replicate the provided inflow profile, it still takes some distance
downstream before the unsteady inflow develops into realistic turbulence.


%###############################################################################
%###############################################################################
%###############################################################################
%###############################################################################
%###############################################################################
%
%\begin{figure}[!hp]
%  \centering
%    \insertimage[width=0.95\textwidth]{sem/sem_box_with_profile}
%    \caption[Creating the SEM box on the inflow plane.]{A virtual box of height
%    \hSEM is constructed such that it straddles the inflow plane starting at the
%    wall. This is shown as the green box in this figure.\\[\capParSep]
%    %
%    During the pre-processor, the wall distances for all boundary points on the
%    SEM boundary are computed. These wall distances are sorted to separate those
%    $\leq\hSEM$ from those $>\hSEM$.  A new subsonic inflow boundary condition
%    is created from the flow conditions specified by \inptext{BC\_Input} and
%    all boundary points outside the SEM profile are moved to this new boundary
%    condition. This ensures that 1) the minimum number of boundary points are
%    contained within the SEM boundary condition, and 2) the subsonic inflow
%    outside the SEM is a not overprescribed allowing for characteristic
%    information to both enter and leave the computational domain through the
%    boundary.\\[\capParSep]
%    %
%    It is best to make sure that the flow conditions specified in
%    \inptext{BC\_Input} for an SEM boundary condition are consistent with the
%    flow conditions at the top of the SEM profile. Stability issues can occur
%    if there is a mismatch between the SEM profile and the corresponding
%    subsonic inflow boundary condition.\\[\capParSep]
%    %
%    The profile data in the SEM profile file is not required to only contain the
%    boundary layer part of the flow. The profile is allowed to contain as much
%    of the main freestream flow as desired. There can be some inefficiencies
%    from including some of the freestream flow in the SEM profile. The height of
%    the SEM box is based on the height of the SEM profile, \hSEM; and the number
%    of synthetic eddies is based on the volume of the SEM box, \vSEM; therefore,
%    including parts of the freestream flow will increase the number of synthetic
%    eddies in the SEM box. In the freestream, the turbulent stresses are
%    generally near-zero.  When the SEM algorithm computes a random fluctuation
%    for a boundary point in the freestream section, the correlation converting
%    this to a physical fluctuation will result in a non-fluctuating velocity
%    profile since the turbulence is zero. All of this is to say that the
%    freestream parts of the SEM profile will waste computational resources
%    trying to compute velocity fluctuations where there are no fluctuations.}
%    \label{FIG:SEM:AddedBox}
%\end{figure}
%
%%-------------------------------------------------------------------------------
%
%\begin{figure}[!hp]
%  \setlength{\imageheight}{0.35\textheight}
%  \centering
%  \insertimage[height=\imageheight]{sem/sem_all_eddies}
%  \captionsetup{singlelinecheck=off}
%  \caption[All eddies and their individual zones of influence.]{This figure is
%  illustrating a group of eddies (red and blue curly arrows) in the vicinity of
%  a face point (large black dot) on an SEM inflow plane. The view is looking
%  downstream with the axial flow, so each eddy is moving into the page. Each of
%  the large black circles is centered around one of the eddies, and the circle
%  represents the ``zone of influence'' for the eddy.\\[\capParSep]
%  %
%  Each eddy has a ``zone of influence'' distance based on the input \lSEM that
%  extends equally in each coordinate direction.}
%  \label{FIG:SEM:ZoI:All}
%\end{figure}
%
%%-------------------------------------------------------------------------------
%
%\begin{figure}[!hp]
%  \setlength{\mpheight}{0.49\textheight}
%  \setlength{\mpwidth}{1.0\textwidth}
%  \setlength{\imagewidth}{0.45\textwidth}
%  \centering
%  \begin{minipage}[c][\mpheight]{\mpwidth}
%    \centering
%    \null
%    \subfig{\imagewidth}{
%      \insertimage[width=\imagewidth]{sem/sem_my_eddies_0}
%      \caption{Eddies that will contribute to the SEM fluctuations for this face
%               point.}
%      \label{FIG:SEM:ZoI:Separated:mine0}
%    }
%    \hfill
%    \subfig{\imagewidth}{
%      \insertimage[width=\imagewidth]{sem/sem_ignored_eddies_0}
%      \caption{Eddies that will not contribute to the SEM fluctuations for this
%               face point.}
%      \label{FIG:SEM:ZoI:Separated:igrnored0}
%    }
%    \null
%  \end{minipage}
%  \caption[Separating out eddies that will not influence face point.]{
%    For a given face point, the eddies with the SEM box can be separated into
%    those that will contribute to the SEM fluctuations at the face point and
%    those that do not contribute.\\[\capParSep]
%    %
%    During the preprocessor phase, after the SEM box is created and populated
%    with synthetic eddies, each face point on the SEM inflow plane will loop
%    through the eddies and collect all the eddies whose zone of influence could
%    potentially include the face point at any time as the eddy propagates
%    through the SEM box. Figure (a) shows all of the eddies that will influence
%    the face point, whereas figure (b) shows all of the eddies that will not
%    influence the face point.}
%  \label{FIG:SEM:ZoI:Separated0}
%\end{figure}
%
%%-------------------------------------------------------------------------------
%
%\begin{figure}[!hp]
%  \setlength{\imageheight}{0.35\textheight}
%  \centering
%  \insertimage[height=\imageheight]{sem/zone_of_influence}
%  \caption[Zone of influence point-of-view for face point.]{When computing the
%  velocity fluctuations at each face point due to the eddies influencing a given
%  point, looping through may thousands of eddies to find the dozens or so that
%  could actually influence the face point at a given time becomes
%  computationally expensive across thousands of face points on the SEM
%  boundary.\\[\capParSep]
%  %
%  To improve the efficiency of this operation, a localized linked-list database
%  is created for each face point that contains a list of just those eddies whose
%  zone of influence might affect the respective face point. Since these eddies
%  are constantly being regenerated at a random location after propagating
%  through the SEM box, each of these linked-lists is constantly being updated.
%  By using these localized linked-lists, each face point loops through only
%  those eddies that can have a possible influence on the face point when
%  computing the SEM velocity fluctuations, thus significantly improving the
%  efficiency of the SEM method.}
%  \label{FIG:SEM:ZoI:FacePoint}
%\end{figure}

%-------------------------------------------------------------------------------

%###############################################################################
%TEMP TEMP TEMP TEMP TEMP TEMP TEMP TEMP TEMP TEMP TEMP TEMP TEMP TEMP TEMP TEMP
%TEMP TEMP TEMP TEMP TEMP TEMP TEMP TEMP TEMP TEMP TEMP TEMP TEMP TEMP TEMP TEMP
%TEMP TEMP TEMP TEMP TEMP TEMP TEMP TEMP TEMP TEMP TEMP TEMP TEMP TEMP TEMP TEMP
%TEMP TEMP TEMP TEMP TEMP TEMP TEMP TEMP TEMP TEMP TEMP TEMP TEMP TEMP TEMP TEMP
%TEMP TEMP TEMP TEMP TEMP TEMP TEMP TEMP TEMP TEMP TEMP TEMP TEMP TEMP TEMP TEMP
%###############################################################################
\clearpage
%###############################################################################
%TEMP TEMP TEMP TEMP TEMP TEMP TEMP TEMP TEMP TEMP TEMP TEMP TEMP TEMP TEMP TEMP
%TEMP TEMP TEMP TEMP TEMP TEMP TEMP TEMP TEMP TEMP TEMP TEMP TEMP TEMP TEMP TEMP
%TEMP TEMP TEMP TEMP TEMP TEMP TEMP TEMP TEMP TEMP TEMP TEMP TEMP TEMP TEMP TEMP
%TEMP TEMP TEMP TEMP TEMP TEMP TEMP TEMP TEMP TEMP TEMP TEMP TEMP TEMP TEMP TEMP
%TEMP TEMP TEMP TEMP TEMP TEMP TEMP TEMP TEMP TEMP TEMP TEMP TEMP TEMP TEMP TEMP
%###############################################################################










